% Section 1.1, from lectures/L01/appendix/a_number_systems.md.

\section{Talsystem}
\label{c1:app:a}

\subsection{Digitalt och analogt}
\label{c1:a1}
En \textbf{analog} signal kan anta vilket värde som helst inom ett intervall. Spänningen från en
temperaturgivare kan vara 1,2~V, 1,21~V eller 1,2137~V, och varje liten ändring betyder något. Det
gör också varje liten störning: brus från en motor i närheten flyttar signalen, och mottagaren
kan inte veta vilken del av värdet som är mätning och vilken som är störning.

En \textbf{digital} signal antar bara två värden, som vi kallar \code{0} och \code{1}. I en krets
som matas med 5~V betyder ungefär 0~V en nolla och ungefär 5~V en etta, och allt som ligger
tillräckligt nära någon av dem läses som exakt \code{0} eller exakt \code{1}. En störning på några
tiondels volt ändrar därför ingenting alls.

\textbf{Den tåligheten mot brus är skälet till att nästan all modern elektronik är digital}, från
mikrokontrollern i en kaffebryggare till styrsystemet i en fabrik. Priset är att en digital signal
bara bär en enda uppgift, ja eller nej, tänd eller släckt. För att kunna representera ett tal, en
temperatur eller en bokstav behövs många digitala signaler tillsammans, och ett sätt att tolka
dem. Det är vad talsystemen och koderna i det här kapitlet handlar om.

En digital signal, eller en siffra i ett binärt tal, kallas en \textbf{bit}, av engelskans
\emph{binary digit}.

\subsection{Positionssystem och talbas}
\label{c1:a2}
Det talsystem vi använder till vardags, det \textbf{decimala} talsystemet, har tio siffror, 0
till 9. Antalet siffror kallas talsystemets \textbf{talbas}, så det decimala talsystemet har
talbasen 10.

Det är också ett \textbf{positionssystem}: en siffras värde beror på var i talet den står. I talet
25 är femman värd 5 och tvåan värd 20. Lägg till en nolla sist, 250, så flyttar båda ett steg åt
vänster och blir värda tio gånger mer: femman 50 och tvåan 200.

Varje position har alltså en \textbf{vikt}, en potens av talbasen. Positionerna numreras från
höger med början på \textbf{noll}:

\begin{narrowtable}{@{}lccc@{}}
  \thead{Position} & 2 & 1 & 0 \\
  \midrule
  Vikt & $10^2 = 100$ & $10^1 = 10$ & $10^0 = 1$ \\
  Siffra i talet 250 & 2 & 5 & 0 \\
  Värde & 200 & 50 & 0 \\
\end{narrowtable}

\noindent Talets värde är summan av varje siffra gånger vikten för dess position:
\[
  250 = 2 \cdot 10^2 + 5 \cdot 10^1 + 0 \cdot 10^0 = 200 + 50 + 0
\]

Att numreringen börjar på noll är ingen godtycklig vana. Den högra positionen har vikten
$B^0 = 1$ i \textbf{alla} talsystem, oavsett talbas $B$, och det är skälet till att den kallas
position noll.

Samma formel gäller alltså för varje talsystem. En siffra $x_n$ på position $n$ i ett tal med
talbasen $B$ är värd
\[
  x_n \cdot B^n
\]
och talets värde är summan av alla siffrornas värden. Den högra siffran, med lägst vikt, kallas
den \textbf{minst signifikanta} siffran, och den vänstra den \textbf{mest signifikanta}.

När det inte framgår av sammanhanget vilket talsystem ett tal är skrivet i, skrivs talbasen som
ett index efter talet: $25_{10}$ är tjugofem decimalt, $11001_2$ är samma tal binärt. I kod och i
den här bokens text används i stället ett prefix, se \secref{c1:a8}.

\subsection{Det binära talsystemet}
\label{c1:a3}
En digital signal har två tillstånd, så den naturliga talbasen i digitaltekniken är 2. Det
\textbf{binära} talsystemet har bara siffrorna 0 och 1, och varje position är värd dubbelt så
mycket som positionen till höger om den:

\begin{narrowtable}{@{}lcccccccc@{}}
  \thead{Position} & 7 & 6 & 5 & 4 & 3 & 2 & 1 & 0 \\
  \midrule
  Vikt & 128 & 64 & 32 & 16 & 8 & 4 & 2 & 1 \\
\end{narrowtable}

\noindent Den raden med vikter är det viktigaste i hela det här kapitlet. Lär dig den utantill,
från höger: 1, 2, 4, 8, 16, 32, 64, 128.

\subsubsection{Från binärt till decimalt}
För att omvandla ett binärt tal till decimalt summerar du vikterna för de positioner där det står
en etta. Nollorna bidrar inte med något och kan hoppas över.

\textbf{Exempel:} $1011_2$ har ettor på position 3, 1 och 0:
\[
  1011_2 = 1 \cdot 8 + 0 \cdot 4 + 1 \cdot 2 + 1 \cdot 1 = 8 + 2 + 1 = 11_{10}
\]

\textbf{Exempel:} $1010\,1110_2$ har ettor på position 7, 5, 3, 2 och 1:

\bookfigure{l01_bit_weights}{Ett 8-bitars tal med bitnumren ovanför och vikterna under: talet
\code{1010 1110} summeras till 174.}{c1:fig:bit_weights}

\[
  1010\,1110_2 = 128 + 32 + 8 + 4 + 2 = 174_{10}
\]

Binära tal skrivs ofta i grupper om fyra siffror, med ett mellanrum emellan, på samma sätt som man
skriver 1\,000\,000 i stället för 1000000. Mellanrummet betyder ingenting; det gör bara talet
lättare att läsa.

\subsubsection{Bitar, nibblar och byte}
Några namn på grupper av bitar som återkommer genom hela boken:

\begin{narrowtable}{@{}lcl@{}}
  \thead{Namn} & \thead{Antal bitar} & \thead{Exempel} \\
  \midrule
  bit & 1 & \code{1} \\
  nibble & 4 & \code{1011} \\
  byte & 8 & \code{1010 1110} \\
  ord (på AVR-processorn) & 16 & \code{0000 0011 1110 1000} \\
\end{narrowtable}

\noindent Den vänstra biten, med störst vikt, kallas \textbf{MSB} (\emph{most significant bit}).
Den högra, med vikten 1, kallas \textbf{LSB} (\emph{least significant bit}). I en byte är MSB
bit~7 och LSB bit~0.

Två saker går att se direkt på LSB och MSB:
\begin{itemize}
  \item \textbf{LSB avgör om talet är udda eller jämnt.} Alla andra vikter är jämna tal, så bara
    LSB kan göra summan udda.
  \item \textbf{MSB avgör om talet är minst hälften av det största talet.} I en byte är MSB värd
    128, och resten tillsammans som mest 127.
\end{itemize}

\subsubsection{Talområden}
Med $n$ bitar finns det $2^n$ olika kombinationer av ettor och nollor, eftersom varje bit
fördubblar antalet möjligheter. Den minsta är noll, så den största är en mindre än antalet
kombinationer:
\[
  \text{antal värden} = 2^n \qquad \text{största värde} = 2^n - 1
\]

\begin{narrowtable}{@{}ccc@{}}
  \thead{Antal bitar} & \thead{Antal värden} & \thead{Talområde} \\
  \midrule
  1 & 2 & 0--1 \\
  4 & 16 & 0--15 \\
  8 & 256 & 0--255 \\
  10 & 1024 & 0--1023 \\
  16 & 65\,536 & 0--65\,535 \\
\end{narrowtable}

\noindent Du kan kontrollera den största: åtta ettor är $128 + 64 + 32 + 16 + 8 + 4 + 2 + 1 =
255$, vilket är $2^8 - 1$.

Det här är ingen teoretisk detalj. Mikrodatorn i den här boken, ATmega328P, räknar med 8 bitar åt
gången, så ett av dess register kan bara hålla talen 0--255. Vad som händer när ett tal inte ryms
är ämnet för \secref{c1:a7}, och för en stor del av kapitel~\ref{c8:ch}.

\subsection{Från decimalt till binärt}
\label{c1:a4}
Åt andra hållet finns två metoder. Båda ger samma svar; välj den du tycker är enklast.

\subsubsection{Metod 1: vikttabellen}
Skriv upp vikterna, och ta för varje vikt från vänster ställning till om den ryms i det som är kvar
av talet. Ryms den skriver du en etta och drar bort vikten; annars skriver du en nolla.

\textbf{Exempel:} omvandla $23_{10}$ till binärt.

\begin{narrowtable}{@{}lccccc@{}}
  \thead{Vikt} & 16 & 8 & 4 & 2 & 1 \\
  \midrule
  Ryms den i det som är kvar? & 23: ja & 7: nej & 7: ja & 3: ja & 1: ja \\
  Kvar efteråt & 7 & 7 & 3 & 1 & 0 \\
  Bit & 1 & 0 & 1 & 1 & 1 \\
\end{narrowtable}

\noindent Alltså är $23_{10} = 10111_2$. Kontrollera baklänges: $16 + 4 + 2 + 1 = 23$.

\subsubsection{Metod 2: division med 2}
Dela talet med 2 och skriv upp resten, som alltid är 0 eller 1. Dela sedan kvoten med 2 igen, och
fortsätt tills kvoten är 0. Resterna är bitarna, med \textbf{den första resten som LSB}.

\textbf{Exempel:} omvandla $23_{10}$ till binärt igen.

\begin{narrowtable}{@{}ccl@{}}
  \thead{Division} & \thead{Kvot} & \thead{Rest} \\
  \midrule
  23 / 2 & 11 & 1 (bit 0, LSB) \\
  11 / 2 & 5 & 1 (bit 1) \\
  5 / 2 & 2 & 1 (bit 2) \\
  2 / 2 & 1 & 0 (bit 3) \\
  1 / 2 & 0 & 1 (bit 4, MSB) \\
\end{narrowtable}

\noindent Resterna läses \textbf{nedifrån och upp}: $10111_2$, samma svar som med tabellen.

Det vanligaste felet med den här metoden är att läsa resterna uppifrån och ned, vilket ger talet
baklänges, $11101_2 = 29$. Kontrollera därför alltid svaret genom att räkna tillbaka till
decimalt.

\subsubsection{Att fylla ut till en hel byte}
$10111_2$ har fem bitar. Ska talet ligga i ett 8-bitars register fylls det ut med nollor till
vänster: \code{0001 0111}. Nollor till vänster ändrar inte värdet, precis som 023 är samma sak som
23.

\subsection{Det hexadecimala talsystemet}
\label{c1:a5}
Binära tal blir snabbt långa och svårlästa. Ett 16-bitars tal som \code{0100 1010 1000 1100} är
lätt att skriva av fel, och det är nästan omöjligt att se vilket tal det är. Därför skrivs binära
tal nästan alltid i det \textbf{hexadecimala} talsystemet, med talbasen 16.

Det behövs 16 siffror, så efter 0--9 används bokstäverna A--F:

\begin{narrowtable}{@{}ccc@{\hspace{3em}}ccc@{}}
  \thead{Decimalt} & \thead{Binärt} & \thead{Hexadecimalt} &
  \thead{Decimalt} & \thead{Binärt} & \thead{Hexadecimalt} \\
  \midrule
  0 & \code{0000} & 0 & 8 & \code{1000} & 8 \\
  1 & \code{0001} & 1 & 9 & \code{1001} & 9 \\
  2 & \code{0010} & 2 & 10 & \code{1010} & A \\
  3 & \code{0011} & 3 & 11 & \code{1011} & B \\
  4 & \code{0100} & 4 & 12 & \code{1100} & C \\
  5 & \code{0101} & 5 & 13 & \code{1101} & D \\
  6 & \code{0110} & 6 & 14 & \code{1110} & E \\
  7 & \code{0111} & 7 & 15 & \code{1111} & F \\
\end{narrowtable}

\noindent\textbf{Poängen med just talbasen 16 är att $16 = 2^4$: en hexadecimal siffra motsvarar
exakt fyra bitar}, en nibble. Omvandlingen mellan binärt och hexadecimalt kräver därför ingen
räkning alls, bara tabellen ovan.

\subsubsection{Från binärt till hexadecimalt}
Dela upp det binära talet i grupper om fyra bitar, \textbf{med början från höger}, och byt varje
grupp mot sin hexadecimala siffra.

\textbf{Exempel:} $1011\,0110_2$.

\bookfigure{l01_hex_grouping}{Talet \code{1011 0110} uppdelat i två grupper om fyra bitar: den
vänstra gruppen, \code{1011}, blir B och den högra, \code{0110}, blir 6.}{c1:fig:hex_grouping}

\noindent Alltså är $1011\,0110_2 = \text{B6}_{16}$.

Om den sista gruppen till vänster har färre än fyra bitar fyller du ut den med nollor:
$1\,1110_2 = 0001\,1110_2 = \text{1E}_{16}$. Det är därför grupperingen börjar från höger. Börjar
man från vänster hamnar de utfyllande nollorna mitt i talet, där de ändrar värdet.

\subsubsection{Från hexadecimalt till binärt}
Byt varje hexadecimal siffra mot sina fyra bitar.

\textbf{Exempel:} $\text{4A8C}_{16}$.

\begin{narrowtable}{@{}lcccc@{}}
  \thead{Hexadecimal siffra} & 4 & A & 8 & C \\
  \midrule
  Fyra bitar & \code{0100} & \code{1010} & \code{1000} & \code{1100} \\
\end{narrowtable}

\noindent Alltså är $\text{4A8C}_{16} = 0100\,1010\,1000\,1100_2$. Fyra tecken i stället för
sexton, och omvandlingen tar några sekunder när tabellen väl sitter.

\subsection{Omvandling mellan alla tre}
\label{c1:a6}
Mellan binärt och hexadecimalt är det bara att gruppera, som i \secref{c1:a5}. Mellan decimalt och
de andra två finns två vägar.

\subsubsection{Från hexadecimalt till decimalt}
Använd formeln från \secref{c1:a2} med talbasen 16. Vikterna är $16^0 = 1$, $16^1 = 16$,
$16^2 = 256$ och $16^3 = 4096$.

\textbf{Exempel:} $\text{F6}_{16}$, där F är 15:
\[
  \text{F6}_{16} = 15 \cdot 16 + 6 \cdot 1 = 240 + 6 = 246_{10}
\]

\textbf{Exempel:} $\text{2A5}_{16}$, där A är 10:
\[
  \text{2A5}_{16} = 2 \cdot 256 + 10 \cdot 16 + 5 \cdot 1 = 512 + 160 + 5 = 677_{10}
\]

Den andra vägen går via binärt: $\text{F6}_{16} = 1111\,0110_2 = 128 + 64 + 32 + 16 + 4 + 2 =
246$. Den är ofta lika snabb för tvåsiffriga tal, och bra som kontroll.

\subsubsection{Från decimalt till hexadecimalt}
Också här finns två vägar.

\textbf{Via binärt:} omvandla till binärt enligt \secref{c1:a4}, och gruppera sedan i fyror.
$200_{10}$ blir $1100\,1000_2$, som är $\text{C8}_{16}$.

\textbf{Division med 16:} som divisionsmetoden i \secref{c1:a4}, men med 16 i stället för 2.
Resterna är hexadecimala siffror, och även här är den första resten den minst signifikanta.

\begin{narrowtable}{@{}ccl@{}}
  \thead{Division} & \thead{Kvot} & \thead{Rest} \\
  \midrule
  200 / 16 & 12 & 8 \\
  12 / 16 & 0 & 12 = C \\
\end{narrowtable}

\noindent Nedifrån och upp: $\text{C8}_{16}$. Kontroll: $12 \cdot 16 + 8 = 192 + 8 = 200$.

\subsubsection{En sammanfattning}

\begin{booktable}{@{}p{7.2em}LLL@{}}
  \thead{Från \textbackslash{} Till} & \thead{Decimalt} & \thead{Binärt} & \thead{Hexadecimalt} \\
  \midrule
  \textbf{Decimalt} & -- & vikttabellen eller division med 2 & via binärt, eller division med
    16 \\
  \textbf{Binärt} & summera vikterna & -- & gruppera i fyror från höger \\
  \textbf{Hexadecimalt} & siffra gånger $16^n$ & fyra bitar per siffra & -- \\
\end{booktable}

\subsection{Binär addition}
\label{c1:a7}
Binära tal adderas med samma uppställning som decimala: kolumn för kolumn från höger, med en
\textbf{minnessiffra} (en \emph{carry}) som förs vidare till nästa kolumn när summan inte ryms i
en siffra.

Skillnaden är bara att det \enquote{blir tio} redan vid två. Fyra regler räcker:

\begin{narrowtable}{@{}lcc@{}}
  \thead{Summa} & \thead{Resultat} & \thead{Minnessiffra} \\
  \midrule
  0 + 0 & 0 & 0 \\
  0 + 1 & 1 & 0 \\
  1 + 1 & 0 & 1 \\
  1 + 1 + 1 (med minnessiffra) & 1 & 1 \\
\end{narrowtable}

\noindent Den tredje raden är den som känns ovan: $1 + 1 = 10_2$, alltså två, skrivet som en nolla
med en etta i minnet.

\textbf{Exempel:} $90 + 55$, med båda talen som 8-bitars binära tal.

\begin{textdrawing}
  minne:   1111 1100
           0101 1010     (90)
         + 0011 0111     (55)
         -----------
           1001 0001     (145)
\end{textdrawing}

\noindent Kontrollera: $1001\,0001_2 = 128 + 16 + 1 = 145$, och $90 + 55 = 145$.

\subsubsection{När summan inte ryms}
Ett 8-bitars register rymmer 0--255. Vad händer med $240 + 17 = 257$?

\begin{textdrawing}
  minne: 1 1110 0000
           1111 0000     (240)
         + 0001 0001     (17)
         -----------
         1 0000 0001
\end{textdrawing}

\noindent Summan behöver nio bitar. I ett 8-bitars register finns bara plats för de åtta högra,
\code{0000 0001}, alltså 1, och den nionde biten, minnessiffran ut ur bit~7, hamnar utanför.
Resultatet blir $257 - 256 = 1$.

\textbf{Det är inte ett fel i datorn; det är vad åtta bitar betyder.} Talet slår runt, som
vägmätaren i en gammal bil som går från 99\,999 till 00\,000. Mikrodatorn sparar den nionde biten
i en särskild flagga, \emph{carry}, så att programmet kan upptäcka att det hände. Det är ämnet för
kapitel~\ref{c8:ch}, där samma sak kallas \textbf{aritmetisk rundgång}.

\subsubsection{Hexadecimal addition}
Hexadecimala tal kan adderas direkt, med regeln att det \enquote{blir tio} vid 16. Oftast är det
lika enkelt att gå via binärt eller decimalt.

\textbf{Exempel:} $\text{3C}_{16} + \text{0A}_{16}$. Entalskolumnen: C + A = 12 + 10 = 22 = 16 + 6,
så siffran blir 6 med en etta i minnet. Sextontalskolumnen: 3 + 0 + 1 = 4. Summan är
$\text{46}_{16}$, och kontrollen i decimalt är $60 + 10 = 70 = 4 \cdot 16 + 6$.

\subsection{Skrivsätt}
\label{c1:a8}
Samma tal kan skrivas på flera sätt, och vilket som används beror på sammanhanget.

\begin{booktable}{@{}p{8.6em}p{10.4em}L@{}}
  \thead{Skrivsätt} & \thead{Exempel (tjugofem)} & \thead{Används} \\
  \midrule
  Index efter talet & $25_{10}$, $11001_2$, $19_{16}$ & i matematik och i läroböcker \\
  Prefix \code{0b} & \code{0b00011001} & i kod, för binära tal \\
  Prefix \code{0x} & \code{0x19} & i kod, för hexadecimala tal \\
  Prefix \code{$} & \code{$19} & i äldre AVR-assembler, för hexadecimala tal \\
  Utan prefix & \code{25} & alltid decimalt \\
\end{booktable}

Den här boken använder prefixen \code{0b} och \code{0x} i löptext och i kod, eftersom det är så
talen skrivs i assemblerprogrammen från kapitel~\ref{c7:ch} och framåt. Indexen används bara i
härledningar som den här, där båda skrivsätten behövs sida vid sida.

Två saker är värda att notera:
\begin{itemize}
  \item \code{0x19} och \code{19} är olika tal: tjugofem respektive nitton. Ett tal utan prefix är
    decimalt.
  \item Nollor till vänster ändrar inte värdet: \code{0x09} och \code{0x9} är samma tal. I kod
    skriver man ofta ut dem ändå, så att det syns hur många bitar talet är tänkt att ha.
\end{itemize}

\subsubsection{Att kontrollera med kalkylatorn}
Kalkylatorn i Windows har ett läge som heter \textbf{Programmerare}. Där visas samma tal samtidigt
som \code{HEX}, \code{DEC}, \code{OCT} (oktalt, talbas 8, som inte används i boken) och \code{BIN}.
Knappen där det står \code{QWORD}, under talen, väljer hur många bitar kalkylatorn räknar med;
varje klick byter till nästa: \code{QWORD} (64), \code{DWORD} (32), \code{WORD} (16) och
\code{BYTE} (8).

Använd den för att kontrollera dina omvandlingar, aldrig för att ta fram dem: på provet finns ingen
kalkylator, och en omvandling du inte kan göra för hand kan du inte heller felsöka i ett program.
